% TOP_PROBLEM: {"id": 30007257, "title": "Integral-point density on curve character varieties", "category_id": 1, "category": {"id": 1, "name": "number_theory", "display_name": "Number Theory", "description": "Properties of integers, prime numbers, Diophantine equations.", "slug": "number-theory", "order_index": 1, "created_at": "2026-07-31T15:26:25.670Z"}, "status": "open"}
% FIRST_POSTED: 2026-10-06
% ATTEMPT_STATUS: unresolved
% OPENAI_CLAIM: {"schema_version":1,"source_repo":"https://github.com/openai/math","source_commit":"adc7f1241b42e322a6451854ab7e4b4c146bf78a","release_date":"2026-10-06","match":"partial","resolution":"partial","independently_reviewed":false,"second_pass":"n/a","families":[{"family":"027","title":"Potential integral density on curve character varieties","subject":"Number theory","manuscripts":[{"title":"Integral points on character varieties of curves","pdf_path":"preprints/Integral-points-on-character-varieties-of-curves-September-25-2026/paper.pdf","readme_path":"preprints/Integral-points-on-character-varieties-of-curves-September-25-2026/README.md","date":"2026-09-25","theorem_ref":"thm:main, eq:main-field (sections/introduction.tex:35)","statement_sources":[{"path":"preprints/Integral-points-on-character-varieties-of-curves-September-25-2026/build/sections/introduction.tex","line":35,"label":"thm:main, eq:main-field"}]}],"lean":null,"reasoning_trace":null}]}
\documentclass[11pt]{article}
\usepackage[margin=25mm]{geometry}
\usepackage{fontspec}
\setmainfont{Latin Modern Roman}
\usepackage{amsmath,amssymb,mathtools}
\usepackage{unicode-math}
\setmathfont{Latin Modern Math}
\usepackage{enumitem}
\usepackage{xurl,hyperref,bookmark}
\hypersetup{colorlinks=true,urlcolor=blue}
\setcounter{secnumdepth}{0}
\setlength{\parindent}{0pt}
\setlength{\parskip}{5pt}
\setlength{\emergencystretch}{3em}
\begin{document}
\section*{Integral-point density on curve character varieties}

\subsection{OpenAI's claim}
OpenAI presents the following claims relevant to this question, with manuscript numbers following paper order in Sources. In manuscript 1 (family 027; 2026-09-25), the asserted result is that every component of the $\mathrm{SL}_r$ character variety of a smooth complex curve has Zariski-dense integral points after one finite field extension, including prescribed exact quasi-unipotent boundary classes.

\subsection{Mathematical statement of the claimed result}
\paragraph{Manuscript 1.}
Take $r\ge1$ and an algebraic curve $X/\mathbb C$ that is connected and smooth. Let $M_r(X)$ be the natural integral $\operatorname{SL}_r$ character scheme. At the punctures of a smooth projective completion, prescribe quasi-unipotent conjugacy classes $C_i\subset\operatorname{SL}_r(\mathbb C)$ and let $Y$ be their reduced relative locus of semisimple representations. Choose a number field $K\subset\mathbb C$ containing all boundary eigenvalues. The claim gives one finite number-field extension $L/K$ such that
\[
 \overline{M_r(X)(\mathcal O_L)\cap Y(\mathbb C)}^{\rm Zar}=Y.
\]
Integrality means extension to the ambient character scheme over the full integer ring, not avoidance of smaller Jordan strata at all reductions. Prescribing only boundary characteristic polynomials with root-of-unity roots, or no boundary data, also works. For the relative cases and for projective $X$, distinct primes $p,q>r$ permit
\[
 F=K(\sqrt2,\zeta_{pq}),\qquad
 L=F(\zeta_r,\{u^{1/r}:u\in\mathcal O_F^{\times}\}).
\]
Without boundary data one may start with $K=\mathbb Q$.

\subsection{Outline of the argument}
For manuscript 1 (family 027; 2026-09-25), filtered local systems are cut along seams and glued with integral splittings. Rank induction handles closed and nonclosed facets together, while flag comparisons recover exact nonsemisimple boundary data rather than only its semisimplification.

\subsection{Scope relative to this problem}
The construction retains the prescribed quasi-unipotent boundary classes and obtains integral density after one finite field extension. This catalogue does not establish complete coverage of the full character-variety statement and records a partial claim.

\subsection{Formal verification}
No Lean formalization is listed at this snapshot for manuscript 1 (family 027; 2026-09-25).

This site has not executed the Lean code. Described statement coverage is not an independent proof check.

\subsection{Status caveat}
The manuscripts were produced by an internal OpenAI model. OpenAI cautions that results without formalization can contain mistakes. This site has not independently checked the proofs; the label records the scope of OpenAI's claim.

\subsection{Sources}
\begin{itemize}
\item OpenAI, \emph{Integral points on character varieties of curves}, 2026-09-25; citation key \texttt{OAI:Integral-points-on-character-varieties-of-curves-September-25-2026}. \href{https://github.com/openai/math/blob/adc7f1241b42e322a6451854ab7e4b4c146bf78a/preprints/Integral-points-on-character-varieties-of-curves-September-25-2026/paper.pdf}{PDF (GitHub)}; \href{https://github.com/openai/math/raw/adc7f1241b42e322a6451854ab7e4b4c146bf78a/preprints/Integral-points-on-character-varieties-of-curves-September-25-2026/paper.pdf}{PDF (direct)}; \href{https://github.com/openai/math/blob/adc7f1241b42e322a6451854ab7e4b4c146bf78a/preprints/Integral-points-on-character-varieties-of-curves-September-25-2026/README.md}{Manuscript page and citation}; \href{https://github.com/openai/math/blob/main/preprints/Integral-points-on-character-varieties-of-curves-September-25-2026/paper.pdf}{Latest version}; Statement source: \href{https://github.com/openai/math/blob/adc7f1241b42e322a6451854ab7e4b4c146bf78a/preprints/Integral-points-on-character-varieties-of-curves-September-25-2026/build/sections/introduction.tex\#L35}{thm:main, eq:main-field}.
\item \href{https://github.com/openai/math/blob/adc7f1241b42e322a6451854ab7e4b4c146bf78a/CONTENTS.md}{OpenAI catalogue entry}.
\end{itemize}
\end{document}
